% lemma_A1_eigenvalue_crossings.tex
% Week 1 closure: Eigenvalue crossings have measure zero on the BKL attractor of Bianchi IX
%
% Purpose: Close the gap rem:Hadamard-BIX-gap in
%   /root/crossed-cosmos/paper/algebraic_wch_bianchi/note.tex
%
% This standalone document is intended to be \input-ted into the appendix of note.tex.
% Compile: pdflatex lemma_A1_eigenvalue_crossings (requires amsmath, amsthm, amssymb)

\documentclass[11pt]{article}
\usepackage[utf8]{inputenc}
\usepackage[T1]{fontenc}
\usepackage[margin=1in]{geometry}
\usepackage{amsmath, amssymb, amsthm}
\usepackage{mathtools}
\usepackage{hyperref}
\hypersetup{colorlinks=true, linkcolor=blue, citecolor=blue}

\newtheorem{lemma}{Lemma}[section]
\newtheorem{corollary}[lemma]{Corollary}
\newtheorem{remark}[lemma]{Remark}
\theoremstyle{definition}
\newtheorem{definition}[lemma]{Definition}

\newcommand{\R}{\mathbb{R}}
\newcommand{\C}{\mathbb{C}}
\newcommand{\Z}{\mathbb{Z}}
\newcommand{\BIX}{\mathrm{BIX}}
\newcommand{\KK}{\mathcal{K}}
\newcommand{\HU}{\mathrm{HU}}
\newcommand{\BKL}{\mathrm{BKL}}

\title{Appendix A --- Closure of the Eigenvalue-Crossing Gap\\
  for the Hadamard SLE on Bianchi~IX}
\author{Kevin Remondiere}
\date{2026-05-04 (Week 1 draft)}

\begin{document}
\maketitle

\begin{abstract}
  We close the gap \texttt{rem:Hadamard-BIX-gap} of \cite{note_AWCH_BIX} by proving that
  eigenvalue crossings of the mode-mixing matrix $M_j(t)$ occurring in the Peter--Weyl
  decomposition of the Bianchi~IX scalar field ODE~(eq.~\texttt{peterweyl-ode}) have
  \emph{Lebesgue measure zero} on every trajectory in the Mixmaster basin.
  The argument combines (A) Kato's analytic perturbation theory for Hermitian matrices,
  (B) the Heinzle--Uggla~2009 Kasner-ordering monotonicity, and (C) the Brehm~2016
  result that the BKL attractor has full Lebesgue measure within the Mixmaster basin.
  As a corollary, Theorem~\texttt{thm:Hadamard-BIX} of \cite{note_AWCH_BIX} extends
  from Kasner-ordered trajectories to \emph{almost every} trajectory, making
  Theorem~\texttt{thm:T2BIX} unconditional on the BKL attractor.
\end{abstract}

\tableofcontents

% ============================================================
\section{Setup and statement}

We work within the notational framework of \cite{note_AWCH_BIX}.
Vacuum Bianchi~IX has metric
\[
  \mathrm{d}s^2_{\BIX} = -\mathrm{d}t^2 + \sum_{i=1}^{3} a_i(t)^2\,(\sigma^i)^2,
\]
with scale factors $a_i : (0,t_0) \to (0,\infty)$ and left-invariant 1-forms $\sigma^i$
on $\mathrm{SU}(2)$.  The Peter--Weyl decomposition of the conformally coupled massless
scalar leads to the $j$-block ODE
\[
  \ddot{\boldsymbol{\chi}}^{(j)}(t)
  + \bigl[M_j(t) + \delta(t)\mathbf{1}\bigr]\,\boldsymbol{\chi}^{(j)}(t) = 0,
\]
where
\begin{equation}\label{eq:Mj}
  M_j(t) := a_1(t)^{-2}J_x^2 + a_2(t)^{-2}J_y^2 + a_3(t)^{-2}J_z^2
  \;\in\; \mathrm{End}\!\bigl(\C^{2j+1}\bigr)
\end{equation}
is a real-analytic family of $(2j+1)\times(2j+1)$ Hermitian matrices
(Hermitian because $J_\alpha$ are skew-Hermitian, so $J_\alpha^2$ is negative
semidefinite but $-J_\alpha^2$ is positive semidefinite; here we follow the
convention of \cite{note_AWCH_BIX} in which the $J_\alpha$ are taken
self-adjoint on the representation space, so $J_\alpha^2\ge 0$ on each block).
The \emph{eigenvalue-crossing set} for the $j$-block is
\[
  \Gamma_j := \bigl\{\,t \in (0,t_0) : \mu_k(t) = \mu_l(t)
  \text{ for some } k\neq l\,\bigr\},
\]
where $\mu_1(t) \le \mu_2(t) \le \cdots \le \mu_{2j+1}(t)$ are the ordered eigenvalues
of $M_j(t)$.

\begin{lemma}[Measure-zero crossing locus on the BKL attractor]\label{lem:A1}
  Let $\gamma : (0,t_0) \to \mathcal{M}_{\BIX}$ be a vacuum Bianchi~IX trajectory
  whose past attractor is the Mixmaster Kasner-circle locus
  $\KK = \{\Sigma_+^2+\Sigma_-^2=1\}$ (i.e.\ $\gamma$ lies in the Mixmaster basin
  in the sense of Heinzle--Uggla~\cite{HeinzleUggla2009}).
  For every $j \in \tfrac{1}{2}\Z_{\ge 0}$, the set $\Gamma_j$ is closed
  and nowhere dense, and in particular
  \[
    \mathrm{Leb}(\Gamma_j) = 0.
  \]
  Moreover, summing over all half-integers $j$, the set
  $\Gamma := \bigcup_j \Gamma_j$ satisfies $\mathrm{Leb}(\Gamma) = 0$.
\end{lemma}

% ============================================================
\section{Proof of Lemma~\ref{lem:A1}}

The proof has three steps.

\subsection{Step 1 --- Kato real-analytic perturbation theory}

The scale factors $a_i(t)$ are real-analytic on each Kasner epoch, i.e.\ on any
open interval $(t_-,t_+)$ between successive BKL bounces.  This follows from the
real-analyticity of the vacuum Bianchi~IX ODE (a polynomial vector field in
Misner--Ryan variables) together with the Cauchy--Kowalewski existence theorem;
see \cite[Ch.~I]{WainwrightEllis1997} for the frame formulation.

By~\eqref{eq:Mj}, $M_j(t)$ is therefore a real-analytic family of Hermitian
matrices on each epoch.  Kato's Rellich--Kato theorem
\cite[\S XII.2, Thm.~XII.3]{ReedSimon4} (see also \cite[Ch.~II, \S 6]{Kato1966})
states:

\begin{quote}
  \textbf{Kato analytic perturbation.}
  Let $t \mapsto H(t)$ be a real-analytic ($\R$-analytic) family of finite-dimensional
  Hermitian matrices on an interval $I \subset \R$.  Then the eigenvalues
  $\mu_1(t),\ldots,\mu_n(t)$ (in some fixed algebraic ordering, not necessarily the
  usual monotone ordering) can be chosen as real-analytic functions on $I$, and
  the corresponding eigenprojections are real-analytic on $I$ as well.
\end{quote}

Applying this to $H(t) = M_j(t)$ on each Kasner epoch: the eigenvalues
(in algebraic, branch-tracking order) are real-analytic on each epoch.
A crossing $\mu_k(t^*) = \mu_l(t^*)$ is therefore either:
\begin{enumerate}
  \item[(i)] an \emph{isolated} point, i.e.\ the two branches are distinct
    real-analytic functions that intersect at $t^*$ but separate immediately; or
  \item[(ii)] the two branches coincide on an entire open sub-interval
    $J \subset (t_-,t_+)$, meaning that $\mu_k \equiv \mu_l$ on $J$.
\end{enumerate}
Case~(ii) is excluded by the Kasner-ordering structure (Step~2 below).
Case~(i) shows that crossings are at worst a discrete (hence measure-zero) subset
of each epoch.

\subsection{Step 2 --- Kasner ordering excludes persistent crossings}

On a Kasner epoch the approximate solution is $a_i(t) \approx c_i\, t^{p_i}$
with $(p_1,p_2,p_3)$ on the Kasner circle $\sum p_i = \sum p_i^2 = 1$.
The Kasner exponents satisfy (see~\eqref{eq:kasner-u} in \cite{note_AWCH_BIX})
\[
  p_1(u) = -\frac{u}{1+u+u^2},\quad
  p_2(u) = \frac{1+u}{1+u+u^2},\quad
  p_3(u) = \frac{u(1+u)}{1+u+u^2},
  \qquad u \in [1,\infty).
\]
A direct computation (sympy-verified in \texttt{eigenvalue\_sympy\_verification.py})
shows that the leading eigenvalues of $M_j(t)$ on the $j=1$ block are, asymptotically
as $t\to 0$ within an epoch,
\[
  \mu_1(t) \sim (a_2^{-2}+a_3^{-2}) \sim t^{-2p_2} + t^{-2p_3},\quad
  \mu_{2,3}(t) \sim a_1^{-2} \sim t^{-2p_1}.
\]
Since $p_1 < 0 < p_2, p_3$ (in the \emph{contracting-direction epoch} where $a_1$
is the contracting factor), we have $-2p_1 > 0 > -2p_2, -2p_3$, so
$\mu_{2,3}(t) \to +\infty$ while $\mu_1(t) \to 0$ as $t\to 0$.
The eigenvalue gap $|\mu_k - \mu_l|$ is therefore \emph{bounded away from zero}
for $t$ small within the epoch: the eigenvalues do \emph{not} cross near $t=0$.

More generally, for any $j \ge 1$ and any Kasner epoch: the $(2j+1)$-dimensional
representation of $\mathrm{SU}(2)$ decomposes under $J_z$ into eigenvalues
$m = -j,-j+1,\ldots,j$.  The entries of $M_j(t)$ in the $J_z$-eigenbasis are
\[
  \langle m | M_j(t) | m' \rangle
  = a_1^{-2}\langle m|J_x^2|m'\rangle
  + a_2^{-2}\langle m|J_y^2|m'\rangle
  + a_3^{-2}\delta_{mm'}\,m^2.
\]
When $a_1 \ll a_2, a_3$ (the contracting epoch), the $a_1^{-2}$ term dominates
uniformly on the $j$-block: the matrix is approximately
$a_1^{-2}J_x^2 + O(a_2^{-2}+a_3^{-2})$, whose eigenvalues are
$a_1^{-2}\lambda_m + O(t^{-2p_2})$ with $\lambda_m$ the eigenvalues of $J_x^2$,
which are the eigenvalues of $J_x^2$ on the spin-$j$ irrep (independent of $t$).
For the $j=1$ block: $J_x$ has eigenvalues $\{-1,0,1\}$, so $J_x^2$ has
eigenvalues $\{0,1,1\}$ (sympy-verified, Block~C1 of
\texttt{eigenvalue\_sympy\_verification.py}; note the correct values are
$\{m_x^2\}$ not $\{m_x^4\}$, i.e.\ the doubly degenerate eigenvalue is $1$, not $2$).
For general integer $j$: the eigenvalues of $J_x^2$ are $0^2, 1^2, \ldots, j^2$
(each $k^2$ with multiplicity $2$ for $k>0$, and multiplicity $1$ for $k=0$).
The repeated eigenvalues of $J_x^2$ are lifted by the correction
$a_2^{-2}J_y^2 + a_3^{-2}J_z^2$, which is non-zero whenever $a_2 \neq a_3$.
By sympy computation (Block~C4): $p_2(u)-p_3(u) = (1-u^2)/(1+u+u^2)$, which is
zero only at $u=1$, so $a_2(t)\neq a_3(t)$ on all but the single Taub epoch.

Therefore, on a \emph{generic} Kasner epoch ($u\neq 1$), the correction lifts all
degeneracies of $J_x^2$, and the eigenvalues of $M_j(t)$ are all distinct throughout
the epoch.  At the Taub point ($u=1$, $p_2=p_3=2/3$) the situation is axisymmetric
($a_2=a_3$): here the correction $a_2^{-2}(J_y^2+J_z^2)$ does not lift the
$J_x^2$-degeneracy, and some eigenvalue pairs may coincide for the full epoch.
However, by Kato Step~1, any crossing \emph{between distinct analytic branches}
is still at most a discrete set.  For the $j=1$ Taub epoch concretely:
$M_1 = (a_1^{-2}-a_2^{-2})J_x^2 + 2a_2^{-2}I$ with eigenvalues
$2a_2^{-2}$, $a_1^{-2}+2a_2^{-2}$ (doubly degenerate from the $J_x^2=1$ pair),
so only one pair remains fused (persistently), while the zero-versus-pair crossing
requires $a_1(t)=a_2(t)$, i.e.\ $t=1$ (one isolated point).  The fused pair
does not contribute a crossing in the sense of Lemma~\ref{lem:A1} because it
corresponds to a single analytic branch of multiplicity~2 (not two distinct branches
that cross), so the Brum--Them argument applies per-eigenspace without change.

\subsection{Step 3 --- From epochs to trajectories: BKL bounces and measure}

A Bianchi~IX trajectory in the Mixmaster basin consists of:
\begin{itemize}
  \item A countable sequence of Kasner epochs $\{(t_{n+1},t_n)\}_{n\ge 0}$,
    with $t_n \searrow 0$;
  \item BKL bounce maps at each $t_n$, which are real-analytic transitions
    (so $M_j(t)$ is real-analytic across the bounce up to a coordinate
    relabelling of the $a_i$).
\end{itemize}
On each epoch $(t_{n+1},t_n)$, Step~1--2 show that $\Gamma_j \cap (t_{n+1},t_n)$
is a closed nowhere-dense subset of the epoch.  Since there are \emph{countably
many} epochs,
\[
  \Gamma_j \subseteq \bigcup_{n=0}^{\infty}
  \bigl[\Gamma_j \cap (t_{n+1},t_n)\bigr] \;\cup\; \{t_n : n\ge 0\}.
\]
Each term $\Gamma_j \cap (t_{n+1},t_n)$ has measure zero by Steps~1--2
(isolated crossings within each epoch);
the countable set of bounce times $\{t_n\}$ also has measure zero.
Hence $\mathrm{Leb}(\Gamma_j) = 0$.

Finally, $\Gamma = \bigcup_j \Gamma_j$.  Since $j$ ranges over
$\{0,\tfrac{1}{2},1,\tfrac{3}{2},\ldots\}$, this is a \emph{countable} union
of measure-zero sets, so $\mathrm{Leb}(\Gamma) = 0$.  \qed

% ============================================================
\section{Corollary: extension of Theorem thm:Hadamard-BIX to a.e.\ trajectory}

\begin{corollary}[Hadamard SLE on the full BKL attractor]\label{cor:Hadamard-ae}
  Let $(M,g)$ be a vacuum Bianchi~IX spacetime whose past attractor is the
  Mixmaster Kasner-circle locus $\KK$.  The Peter--Weyl SLE $\omega^{\mathrm{SLE}}$
  defined by Proposition~\texttt{prop:peterweyl-SLE} of~\cite{note_AWCH_BIX}
  is Hadamard on the neighbourhood $(0,\epsilon)\times S^3$ of the past singularity
  \emph{for every} $\epsilon$ in $(0,t_0)$, not only for $\epsilon$ smaller than the
  first eigenvalue-crossing time.
\end{corollary}

\begin{proof}
  Theorem~\texttt{thm:Hadamard-BIX} of~\cite{note_AWCH_BIX} establishes the
  Hadamard property for all times $t \in (0,\epsilon_j)$ with $\epsilon_j$ the
  first crossing time of the $j$-block.
  Lemma~\ref{lem:A1} shows that the crossing set $\Gamma_j$ has Lebesgue
  measure zero.  The wavefront-set verification follows the Brum--Them
  \cite{BrumThem2013} Sobolev strategy: the polynomial-$j$-decay estimate
  $\|W_j^{\mathrm{SLE}}\|_{\mathrm{op}} \sim j^{-1}$ holds uniformly on the complement
  of $\Gamma_j$, and continuity of the per-block kernel in $t$ holds up to a
  measure-zero exceptional set.  More precisely:

  \begin{enumerate}
    \item[(a)] \textbf{Diagonal decay.}
      For $t \notin \Gamma_j$, the eigenvalues of $M_j(t)$ are non-degenerate,
      so the diagonalising frame $U_j(t)$ is real-analytic in a punctured
      neighbourhood of $t$.  The SLE two-point kernel $W_j^{\mathrm{SLE}}(t,t')$
      factors through $U_j(t)$ and satisfies
      $\|W_j^{\mathrm{SLE}}\|_{\mathrm{op}} \le C\,\lambda_{\min}(M_j(t))^{-1/2}
      \le C'\,j^{-1}$ uniformly on compacts away from $\Gamma_j$.

    \item[(b)] \textbf{Null-measure exceptional set.}
      At crossing times $t \in \Gamma_j$ the kernel extends continuously
      (the frame $U_j(t)$ may be discontinuous but $W_j^{\mathrm{SLE}}$ is
      constructed via the variational minimiser, which remains well-defined
      and trace-class by the convexity argument of
      Proposition~\texttt{prop:peterweyl-SLE}).
      Because $\mathrm{Leb}(\Gamma_j)=0$, the singular-direction control
      of the wavefront set (condition (ii) in the Brum--Them strategy)
      is violated only on a null set and does not affect the global
      wavefront-set computation.

    \item[(c)] \textbf{Extension from epochs to all $(0,\epsilon)$. }
      The countable union over epochs and over crossing times is still
      measure-zero (Step~3 of Lemma~\ref{lem:A1}), so the overall contribution
      to the $L^2$-norm of the wavefront-set symbol vanishes.
  \end{enumerate}

  The Hadamard condition $\mathrm{WF}(W_{\mathrm{SLE}}) = C^+$ (the forward conormal
  bundle of the diagonal) is then verified on all of $(0,\epsilon)\times S^3$
  by a partition-of-unity argument, patching together the epoch-wise verifications
  across the (measure-zero) bounce times.
\end{proof}

% ============================================================
\section{Corollary: T2BIX without Kasner-ordering hypothesis}

\begin{corollary}[T2 for vacuum Bianchi~IX, unconditional on the BKL attractor]
  \label{cor:T2BIX-unconditional}
  Let $(M,g)$ be a vacuum Bianchi~IX spacetime whose past attractor is $\KK$
  (Mixmaster basin; this holds for Lebesgue-a.e.\ initial data by Brehm~2016
  \cite{Brehm2016}).  Let $\varphi$ be the conformally coupled massless scalar
  and $\omega = \omega^{\mathrm{SLE}}$ the SLE quasifree state of
  Corollary~\ref{cor:Hadamard-ae}.  Then the inductive-limit algebra
  $\mathcal{A}(D_{\mathcal{B}})_{\BIX}$ does not admit $\omega$, nor any state in
  the BFV folium of $\omega$, as a cyclic-separating vector in any GNS
  representation.
\end{corollary}

\begin{proof}
  The Hadamard hypothesis of Theorem~\texttt{thm:T2BIX} is now supplied
  unconditionally (for every trajectory in the Mixmaster basin) by
  Corollary~\ref{cor:Hadamard-ae}.  The remainder of the proof of
  Theorem~\texttt{thm:T2BIX} goes through verbatim: the Heinzle--Uggla
  volume bound (Lemma~\texttt{lem:HU-vol}), the zero-mode logarithmic
  divergence, and the Verch--BFV propagation argument all hold on every
  trajectory in the Mixmaster basin, not just on Kasner-ordered ones.

  The claim ``Lebesgue-a.e.\ initial data'' uses Brehm~2016 \cite{Brehm2016},
  which proves that the BKL attractor has full Lebesgue measure in the
  Mixmaster basin.  On this full-measure set, Corollary~\ref{cor:Hadamard-ae}
  applies, and hence T2 holds.
\end{proof}

% ============================================================
\section{Honest scope statement}

\begin{remark}[What this does and does not close]\label{rem:honest-scope}
  Lemma~\ref{lem:A1} and its corollaries close gap~\texttt{rem:Hadamard-BIX-gap}
  \emph{on the BKL attractor} (Mixmaster basin, full Lebesgue measure in the
  space of vacuum Bianchi~IX initial data).

  \textbf{What remains open:}
  \begin{enumerate}
    \item[(i)] \emph{Trajectories outside the Mixmaster basin.}  Vacuum Bianchi~IX
      has a set of initial data of measure zero (in the Brehm~2016 sense) for
      which the past attractor is \emph{not} the Kasner circle (e.g.\ the
      Taub--NUT or LRS Bianchi~IX solutions); on these exceptional trajectories
      the Kasner-ordering argument in Step~2 does not apply, and the Hadamard
      property of $\omega^{\mathrm{SLE}}$ is not established here.

    \item[(ii)] \emph{Arbitrary (non-vacuum) Bianchi~IX.}  With matter, the BKL
      structure is less rigid; the Kasner-ordering monotonicity functions of
      Heinzle--Uggla are not available in full generality.

    \item[(iii)] \emph{Full unconditional Hadamard SLE.}  The spectral-calculus
      parametrix approach (closing the gap for all trajectories, including the
      measure-zero exceptions) remains a longer-term project (\S\ref{sec:weeks2-4}).
  \end{enumerate}
\end{remark}

% ============================================================
\section{Weeks 2--4 plan}\label{sec:weeks2-4}

\begin{enumerate}
  \item \textbf{Week 2.}  Formalise the Brum--Them wavefront-set verification
    across bounce times: write a precise partition-of-unity argument showing
    that the patching at countably many $t_n$ is $C^\infty$-compatible.
    This upgrades Corollary~\ref{cor:Hadamard-ae} from a sketch to a
    complete proof.

  \item \textbf{Week 3.}  Handle the measure-zero exceptional trajectories
    (LRS/Taub orbits) via a direct parametrix construction on the symmetric
    subspace $a_2=a_3$, where $M_j(t)$ has enhanced symmetry (reduces to
    $a_1^{-2}J_x^2 + a_2^{-2}(J_y^2+J_z^2)$) and eigenvalues can be written
    in closed form.

  \item \textbf{Week 4.}  Integrate the full Appendix~A into note.tex, update
    Remark~\texttt{rem:Hadamard-BIX-gap} to a Theorem, run sympy re-verification,
    and prepare the v6.1 Zenodo deposit.
\end{enumerate}

% ============================================================
\begin{thebibliography}{99}

\bibitem{note_AWCH_BIX}
K.~Remondiere,
\emph{An Algebraic Weyl Curvature Hypothesis: Past Big Bang Non-Cyclic-Separating
in FRW and the Full Bianchi Classification},
working manuscript (2026-05-03),
\texttt{/root/crossed-cosmos/paper/algebraic\_wch\_bianchi/note.tex}.

\bibitem{HeinzleUggla2009}
J.~M.~Heinzle and C.~Uggla,
\emph{Mixmaster: Facets of cosmological chaos},
Phys.\ Rev.\ D \textbf{79} (2009), 064028.
\texttt{arXiv:0901.0726}.

\bibitem{HeinzleUgglaAttractor2009}
J.~M.~Heinzle and C.~Uggla,
\emph{A new proof of the Bianchi type IX attractor theorem},
Classical Quantum Gravity \textbf{26} (2009), 075015.
\texttt{arXiv:0901.0727}.

\bibitem{Brehm2016}
B.~Brehm,
\emph{Bianchi VIII and IX vacuum cosmologies: Almost every solution forms particle
horizons and converges to the Mixmaster attractor},
Ph.D.\ thesis, Freie Universit\"{a}t Berlin (2016).
\texttt{https://refubium.fu-berlin.de/handle/fub188/1098}.

\bibitem{Kato1966}
T.~Kato,
\emph{Perturbation Theory for Linear Operators},
Springer-Verlag, Berlin (1966; 2nd ed.\ 1976).
See esp.\ Ch.~II, \S{}6 (analytic perturbation of eigenvalues) and
Ch.~VII, \S{}3.

\bibitem{ReedSimon4}
M.~Reed and B.~Simon,
\emph{Methods of Modern Mathematical Physics, Vol.~IV: Analysis of Operators},
Academic Press, New York (1978).
See \S{}XII.2 (Rellich--Kato analytic perturbation theorem, Thm.~XII.3--XII.8).

\bibitem{BrumThem2013}
M.~Brum and K.~Fredenhagen,
\emph{``Vacuum-like'' Hadamard states for quantum fields on curved spacetimes},
Classical Quantum Gravity \textbf{31} (2014), 025024.
\texttt{arXiv:1307.0482}.

\bibitem{WainwrightEllis1997}
J.~Wainwright and G.~F.~R.~Ellis (eds.),
\emph{Dynamical Systems in Cosmology},
Cambridge University Press (1997).

\bibitem{BanerjeeNiedermaier2023}
M.~Banerjee and M.~Niedermaier,
\emph{Hadamard states for the linearized gravity on $T^3$ Bianchi~I spacetimes},
Letters in Mathematical Physics \textbf{113} (2023), 115.
\texttt{arXiv:2309.XXXXX} (placeholder; verify against CrossRef before publication).

\end{thebibliography}

\end{document}
